\scanpage{13}
% Foliotation archivistique au crayon, en bas à gauche : 12.
de sorte que
\[
(29)\qquad z\overset{\unsure{\mathrm{dfn}|\mathrm{def}}}{=}\unsure{\mathcal Z|Z}\cap\mathbb G'=\text{Centre de }\mathbb G'
\]
et on a une suite exacte \uline{centrale}
\diagram{\begin{tikzcd}[ampersand replacement=\&,column sep=large]
(30)\quad 1 \arrow[r,"{\deleted{z}}"'] \& z \arrow[r] \& Z \arrow[r] \& \Gamma
\end{tikzcd}}
Donc les solutions de notre pb \inserted{\unsure{relatifs|relatives} à $\Gamma_i\subset\Gamma$} correspondent aux scindages partiels de cette suite exacte
\diagram{\begin{tikzcd}[ampersand replacement=\&,column sep=large]
Z \arrow[r] \& \Gamma \\
\& \Gamma_i \arrow[u,hook] \arrow[ul]
\end{tikzcd}}
Soit $\Gamma_0\subset\Gamma$ l'image de $Z$ dans $\Gamma$
\inserted{($\Gamma_0\leftrightarrow k_0$ \unsure{ss-ext.} finie de $\bar k/k$)}
(qui est un ss-groupe ouvert), donc on a une extension centrale
\diagram{\begin{tikzcd}[ampersand replacement=\&,column sep=large]
(31)\quad 1 \arrow[r,no head] \& z \arrow[r,no head] \arrow[d,equal] \& Z \arrow[r] \& \Gamma_0 \arrow[r] \& 1 \\
\& \operatorname{Centre}(\mathbb G') \& \& \&
\end{tikzcd}}
et il s'agit des scindages partiels de cette extension. Si $z=1$
(centre de $\mathbb G'$ trivial) alors il y a unicité pour $\unsure{\mathrm{tt}|\forall}\ \Gamma_i$, et existence ssi $\Gamma_i\subset\Gamma_0$ --- donc une solution \inserted{unique} \unsure{can-}
